\section{Introduction}
\label{sec:introduction}

Private equity (PE) has grown into a major asset class, with more than \$13 trillion in assets under management \citep{McKinsey2024}. 
Yet rigorous asset-pricing tools for private markets remain comparatively underdeveloped. The central difficulty is that private equity investments do not generate reliable traded-return series. 
Instead, the econometrician observes irregular capital calls, distributions, and reported Net Asset Values (NAVs), often over long horizons and under substantial informational frictions. 
As a result, standard return-based asset-pricing methods that are routine in public markets are difficult to apply in private markets, where the relevant pricing restrictions are expressed in discounted cash flows rather than in observed returns.

Semiparametric stochastic discount factor (SDF) methods provide a natural response to this measurement problem. 
The seminal contribution of \citet{DLP12} shows how to estimate abnormal performance and systematic risk exposure from pooled private-equity cash flows by imposing a zero-net-present-value restriction at fund inception. 
\citet{KN16} develop a closely related framework that links private-equity cash flows to public-market benchmarks while emphasizing the role of fund overlap and cross-sectional dependence. 
These approaches have become central tools in empirical work on private assets because they avoid the strong distributional assumptions required by fully parametric models and can be implemented directly on cash-flow data.

At the same time, existing semiparametric approaches leave open an important practical question: how do cash-flow-based SDF estimators behave in the finite samples researchers actually observe, and which economically meaningful objects can those samples identify with stability? 
Private-equity panels are typically sparse, nonlinear, and heavily overlapping across funds and vintages. 
In addition, recent vintages often contain large reported NAV components, so empirical conclusions may reflect continuation values and sample construction as much as underlying risk. 
In such settings, conclusions about alpha, market exposure, or multifactor structure may depend not only on economics, but also on horizon choice, estimator variance, and the treatment of cross-fund dependence.

This paper studies this identification problem in semiparametric cash-flow-based SDF estimation when the econometrician is not restricted to discounting all pooled cash flows back to fund inception. 
Our starting point is simple. 
In typical private-equity samples, inception-only pricing errors can be highly unstable because the estimating equations are nonlinear and the effective cross-sectional sample size is limited. 
A natural response is to average pricing errors over additional compounding dates rather than evaluate all discounted cash flows only at inception. 
But this is not merely a statistical refinement. Once pricing errors are evaluated at horizons beyond inception, the estimator generally targets a different population object. 
Horizon choice therefore creates a tradeoff between finite-sample stabilization and population target shift: longer horizons can improve finite-sample behavior by smoothing the sample criterion, but they also change the underlying population pricing restriction through what we term the Compounding-Horizon Wedge.

The paper's first contribution is to characterize identification in this setting. 
We derive a generalized pooled-cash-flow pricing framework that makes explicit how the choice of discounting date affects the estimand. 
This formulation nests the inception-only criterion of \citet{DLP12} as a special case. 
It also clarifies the relation to \citet{KN16}: their estimator is closely connected to our framework, but it is not obtained as a literal special case because the order of aggregation differs. 
Framing the problem in this way makes explicit a distinction that, to our knowledge, has not been formalized in the semiparametric private-equity SDF literature: a horizon may improve finite-sample performance without preserving the original inception-based pricing restriction.
Accordingly, we treat horizon averaging as a finite-sample regularization device: it can smooth the nonlinear sample criterion and stabilize estimation, but only by introducing the Compounding-Horizon Wedge and thereby shifting the population target.

Our second contribution is methodological in service of that identification problem. 
We formulate the estimator as a nonlinear Least-Mean-Distance (LMD) problem for pooled, non-traded cash flows \citep{PP97}. 
This objective-function representation is useful because it makes the role of horizon choice transparent, accommodates alternative weighting schemes, and yields a clean comparison to existing semiparametric estimators. 
Building on this representation, we develop an increasing-domain asymptotic framework for inference under dependence induced by overlapping funds and adjacent vintages. 
In private-equity applications, cross-sectional units are neither independent nor naturally ordered in calendar time alone; dependence is tied to economic proximity, such as vintage-year distance and overlap in cash-flow histories. 
Our asymptotic results provide a large-sample benchmark for that setting. 
In the empirical application, however, we treat finite-sample simulation evidence and stability diagnostics as the primary basis for interpretation whenever asymptotic approximations appear unreliable.

The paper's central tension has two distinct margins. 
On one margin, inception-only discounting preserves the original zero-net-present-value pricing restriction but can produce unstable estimates in realistic private-equity samples. 
On the other margin, averaging pricing errors over additional compounding dates can materially improve finite-sample performance, but it does so by changing the population pricing restriction. 
Our goal is not to treat horizon choice as a purely mechanical implementation detail, but to characterize this tradeoff explicitly and provide guidance on how researchers should navigate it in empirical work.

In that sense, this is an identification and stability paper for private-market asset pricing. 
It studies how estimator design, especially the discounting horizon, jointly determines the population object being estimated and the finite-sample stability with which it can be recovered. 
In the private-equity panels we study, this framework yields more stable evidence on broad market exposure than on abnormal performance or richer factor decompositions.

We study this identification and stability problem in three steps. 
First, we derive the generalized pooled-cash-flow objective and the corresponding asymptotic framework. 
Second, we use simulations calibrated to private-equity-like cash-flow environments to evaluate how horizon choice affects finite-sample bias and variance. 
Third, we implement the estimator in the Preqin cash-flow database for North American U.S.-dollar buyout (BO) and venture-capital (VC) funds and assess specification stability using both asymptotic inference and $hv$-block cross-validation diagnostics \citep{R00} tailored to overlapping cash-flow dependence.

Our empirical results yield three main messages, all of which we interpret as diagnostics about identification under currently available private-equity cash-flow panels rather than as definitive statements about the population factor structure of private equity. 
First, including recent vintages with large NAV components can mechanically inflate measured market exposure, underscoring the importance of sample construction in private-equity asset pricing. 
Second, within our sample, parsimonious single-factor market specifications applied to seasoned vintages are the most stable among the specifications we consider. 
Third, $\alpha$ terms and additional factors are only weakly identified in the available data and become highly sensitive to horizon choice and finite-sample noise, even when asymptotic inference is supplemented with cross-validation diagnostics. 
Taken together, these results suggest that, in our sample and under our semiparametric design, current cash-flow panels support more stable inference about broad market exposure in seasoned buyout and venture-capital portfolios than about abnormal performance or richer factor decompositions.

These findings matter for more than estimator design. 
A large literature asks whether private-equity funds outperform public benchmarks and how their risks should be measured (comprehensively summarized by \cite{K19,korteweg2023risk}). 
Our results suggest that these questions cannot be separated from the identification limits imposed by sparse, overlapping, and partially NAV-based cash-flow panels. 
In particular, our evidence does not imply that private-equity alpha is zero, nor does it invalidate prior estimates of outperformance as such. 
Rather, it indicates that semiparametric alpha estimates from pooled cash-flow data can be weakly identified and highly sensitive to implementation choices in samples of the kind currently available. 
That qualification is important for interpreting both positive and negative findings in the broader literature on private-equity outperformance.

The paper contributes to the literature on private-market performance measurement and risk adjustment in three ways. 
Relative to \citet{DLP12}, we generalize the inception-only pricing restriction and provide an asymptotic framework that clarifies large-sample inference under overlapping funds. 
Relative to \citet{KN16}, we show that both approaches are built on closely related cash-flow pricing restrictions, but that different aggregation and weighting choices imply different sample objectives and estimators.
More broadly, the paper contributes to the growing literature on asset-pricing inference with illiquid, non-traded, and intermittently valued cash-flow streams, where the key challenge is not only estimation, but also determining which economically meaningful objects available data can identify with stability.

The remainder of the paper proceeds as follows. 
Section~\ref{sec:literature} reviews the related literature. Section~\ref{sec:spatial_sdf_methodology} develops the generalized semiparametric estimator and the corresponding asymptotic inference framework. 
Section~\ref{sec:empirical_application} discusses implementation choices and studies finite-sample behavior in simulation.
Section~\ref{sec:empirical_results_bovc} presents the empirical application to buyout and venture-capital funds. 
Section~\ref{sec:conclusion} concludes.



\section{Related literature}
\label{sec:literature}

For private equity funds, performance evaluation under the SDF framework remains an active research topic \citep{DLP12, KN16, KN24, giommetti2021risk, GN18, boyer2023discount,korteweg2024private} alongside an expanding list of applications across other private asset classes.
Specifically, the \cite{DLP12} estimator is heavily utilized to price human capital contracts, venture capital, real estate and private equity funds \citep{kroencke2013return,kim2015estimation,farrelly2019risk,lee2019new,zhao2020master,TP24}, while the \cite{KN16} approach has found broad application spanning private equity, infrastructure, real estate investment trusts, collateralized loan obligations, impact funds and private debt \citep{giommetti2021risk,andonov2021institutional,riddiough2022private,boyer2023discount,cordell2023clo,erel2024risk,jeffers2024risk,KN24}.
Previously, SDF performance evaluation was already applied for public stock markets by, e.g., \cite{FFJT02} and for hedge fund returns by \cite{LXZ16}.

Within the private equity literature, SDF-based approaches can be broadly classified into parametric and semiparametric methods.
Parametric models, such as those proposed by \cite{FNP12} and \cite{ACGP18}, explicitly specify the return dynamics and the underlying data generating process of private equity investments. 
While these methods can uncover latent factors or deal-level dynamics, they typically rely on strong distributional assumptions (e.g., log-normality) that may not fully capture the irregular and skewed nature of PE cash flows.

Consequently, semiparametric approaches have gained prominence, starting with the seminal cross-sectional GMM framework of \cite{DLP12}. 
They estimate a linear factor model that minimizes the squared pricing errors of PE vintage year portfolios, explicitly evaluating discounted cash flows at fund inception.
In a nutshell, the \cite{DLP12} estimator determines the estimated alpha and beta parameters of a linear factor model by minimizing the squared net present value (NPV) of $n$ portfolio-level cash-flow streams in a least-squares optimization problem:
\begin{equation}
	\label{eq:dlp12_minimization}
 	\left(\hat{\alpha}, \hat{\beta} \right) = 
 	\arg \min_{\alpha,\beta} \sum_{i=1}^n \left[ {NPV}_i(\alpha,\beta) \right]^2
\end{equation}
with $NPV$ for the $i$th portfolio
\[
{NPV}_i(\alpha,\beta) =
\sum_{t \in \left\{ t_{0,i}, t_{0,i} + 1, \dots \right\}}
\frac{\mathrm{CashFlow}_{i,t}}{
	\prod_{s=t_{0,i}+1}^{t}
	\left( 
	1 + r_{f, s} + \alpha + \beta r_{\mathrm{market}, s}
	\right)
} 
\]
with fund inception date $t_{0,i}$, risk-free rate $r_{f, s}$ and excess market return $r_{\mathrm{market}}$.
Here, the \cite{DLP12} approach relies on the identifying assumption that the model is correctly specified, so that $\mathbb{E}[NPV_i(\theta_0)] = 0$.

Building on the SDF-pricing idea, \cite{KN16} introduced the Generalized Public Market Equivalent (GPME) framework, shifting the estimation of the SDF to public market data while evaluating PE cash flows in a second step. 
Importantly, \cite{KN16} were the first to apply a spatial heteroskedasticity and autocorrelation consistent (SHAC) covariance matrix estimator \citep{C99} to account for the economic distance (cash flow overlap) between PE funds, moving beyond simple cross-sectional bootstrapping.

Further complementing the cash-flow pricing dimension, a strand of literature models the informational frictions and smoothing inherent in reported Net Asset Values (NAVs) \citep{BGK19, JLRS20, brown2023nowcasting}. 
Additionally, \cite{HSS23} leverage a similar generalized framework to evaluate private firms using credit market equivalents.
Finally, for comprehensive methodological and empirical overviews of risk and return assessment in private equity, we refer to \cite{K19,korteweg2023risk}, who also discuss the popular \cite{D79} time-series regression that employs lagged returns to correct for stale valuations.

Our paper's LMD framework contributes to this semiparametric literature by providing a loss-function formulation that encompasses the \cite{DLP12} approach and contextualizes the \cite{KN16} time-series analogue while supplying the rigorous increasing-domain asymptotic inference missing in earlier work (drawing on the spatial theory of \cite{PP97}, \cite{KS11}, and \cite{JP12}). 
While \cite{GSW20} highlighted that NPV-based measures are biased relative to per-period returns, we explicitly formalize this discrepancy as the Compounding-Horizon Wedge, that is, a horizon-induced wedge in the pooled pricing restriction, a concept related to recent literature on discount-rate risk in secondary markets \citep{boyer2023discount}.

% Related conceptually, \citet{APR18} study pseudo-true SDFs in conditional traded-asset pricing models, where the pseudo-true parameter minimizes a conditional Hansen-Jagannathan distance and can be characterized through local factors or their mimicking portfolios. That paper is relevant for our framing because it emphasizes that pseudo-true targets in asset pricing are criterion-dependent objects under misspecification rather than automatically structural parameters. Our setting differs, however, in that the horizon-specific pseudo-target \(\theta_H^\dagger\) arises here because the pooled cash-flow criterion itself changes with the discounting horizon; the central issue is target shift induced by post-inception timing terms in non-traded private-equity cash flows, not selection of a best-fitting misspecified conditional SDF in traded-return data.
